% ---
% title: 创新实践期末考试复盘与知识点
% date: 2026-07-02
% author: Crystal-Sky
% tags: 创新实践, 期末考试, 复盘, 知识点
% summary: 《创新实践与展示》2025-2026 学年期末考试复盘，按题型整理问题、点评和对应知识点。
% slug: innovation-practice-final-review
% pdf: assets/mooc-questions.pdf
% ---

\begin{abstract}
这次考试之前，原本以为试卷会比较依赖 MOOC 原题，真正考出来以后才发现全是PPT 里的概念。所以历年的试卷和mooc其实参考价值很小，除了大题以外。
\end{abstract}

\section{一、单选题}

单选题里有不少 MOOC 原题的影子，但也混进了不少 PPT 细节。

\begin{question}{1. 切换到上级目录的命令}
\textbf{点评：} 这一题比较基础，终端里切换到上一级目录使用 \texttt{cd ..}。

\textbf{知识点：} \texttt{cd} 用于切换目录，\texttt{..} 表示当前目录的上级目录。
\end{question}

\begin{question}{2. 图像形态学中的收缩与像素区域变化}
\textbf{点评：} 我没记错的话这题好像是一个错题，前面强调收缩，然后最后来了个像素增多，疑似拼好题，收缩是腐蚀，像素增多是膨胀。

\textbf{知识点：} 腐蚀会让前景区域变小，常用于去除小噪声；膨胀会让前景区域变大，常用于填补空洞或连接断裂区域。
\end{question}

\begin{question}{3. Python 中的与运算}
\textbf{点评：} python的与运算是and，而不是&

\textbf{知识点：} \texttt{and} 是Python逻辑运算符，判断表达式真假；。
\end{question}

\begin{question}{4. LCD1602 英文字符的编码方式}
\textbf{点评：} 这个题ppt里面好像也没有,可能是mooc视频里讲了？不太清楚，有二进制、十六进制、ASCII和国际码，猜一下ASCII

\textbf{知识点：} LCD1602 显示英文字符时通常按字符编码写入显示数据，常见英文字符可对应 ASCII 编码。
\end{question}

\begin{question}{5. 检测物体遮挡的传感器}
\textbf{点评：} 这题直接选光电传感器，高中光电门好吧。

\textbf{知识点：} U 型光电开关由红外发射器和接收器构成，两者位于 U 型槽两侧并形成光轴。当物体经过 U 型槽并阻断光轴时，光电开关会输出检测信号。
\end{question}

\begin{question}{6. OpenCV 色彩阈值化函数}
\textbf{点评：} 如果刷过 MOOC 题并记住了常见 OpenCV 代码，这题应该能直接选 \texttt{cv2.inRange()}。有意思的是，后面大题里的注释似乎也把答案露出来了。

\textbf{知识点：} \texttt{cv2.inRange(hsv, lower, upper)} 用于判断像素是否落在指定颜色范围内，常用于 HSV 空间下的颜色阈值分割。
\end{question}

\begin{question}{7. 人脸检测与人脸识别的区别}
\textbf{点评：} ppt上原话。

\textbf{知识点：} 人脸检测：指对于任意一幅给定的图像/视频帧，采用一定的策略对其进行搜索以确定其中是否含有人脸，如果是则返回一脸的位置、大小和姿态。（有没有人脸）

人脸识别：基于人的脸部特征信息进行身份识别的一种生物识别技术。（是谁的脸） 
\end{question}

\begin{question}{8. 二维码解码操作的库函数}
\textbf{点评：} 巧了，前一天瞅到了这个： \texttt{pyzbar}。

\textbf{知识点：} \texttt{pyzbar} 的核心解码函数是 \texttt{decode()}，可用于识别二维码、条形码等编码图像。
\end{question}

\begin{question}{9. 树莓派的主存卡类型}
\textbf{点评：} 如果看ppt了就知道是 TF 卡，也就是 Micro SD Card。

\textbf{知识点：} 树莓派常通过 Micro SD/TF 卡安装系统并存储数据。
\end{question}

\begin{question}{10. PCF8591 的功能}
\textbf{点评：} PCF8591 文档首页就写得很清楚：它是 8 位 A/D 和 D/A 转换器。

\textbf{知识点：} PCF8591 是带 I2C 接口的 8 位数据采集器件，常见功能包括多路模拟输入、模拟输出和 A/D、D/A 转换。
\end{question}

\section{二、多选题}

多选题是整张卷子的重灾区。最难的地方不是完全不会，而是不确定到底该不该四选，选多了又害怕。

\begin{question}{1. 属于开关量类型的传感器和器件}
\textbf{点评：} 四个选项有三个带开关：轻触开关、触摸开关、倾斜开关，剩下一个不知道选不选，忘了是个啥了。

\textbf{知识点：} 开关量输出通常只有两种状态，例如高/低电平、通/断、0/1。
\end{question}

\begin{question}{2. \texttt{cv2.flip(img, flipcode)} 中 \texttt{flipcode} 的含义}
\textbf{点评：} 这题算运气好，考前刚好扫到过，直接四选

\textbf{知识点：} \texttt{flipcode=0} 表示沿 x 轴翻转；\texttt{flipcode>0} 表示沿 y 轴翻转；\texttt{flipcode<0} 表示 x、y 轴同时翻转。
\end{question}

\begin{question}{3. Haar 分类器的重要组成}
\textbf{点评：} 不太自信考场上，连着两个四选。

\textbf{知识点：} Haar 分类器通常包括 Haar-like 特征、积分图方法、AdaBoost 训练和级联分类器。
\end{question}

\begin{question}{4. 二维码相比一维码的优势}
\textbf{点评：} 这题比较遗憾。当时没理解“二维码纠错能力更强”，以为二维矩阵错一块会更难恢复。实际上二维码正是因为有纠错编码，容错能力更强。

\textbf{知识点：} 一维码主要在一个方向表达信息；二维码在水平和垂直两个方向存储信息，容量更大，可表达更复杂的数据，并且通常具有更强的纠错能力。
\end{question}

\begin{question}{5. Python 循环语句}
\textbf{点评：} 这题比较基础，循环语句就是 \texttt{for} 和 \texttt{while}。\texttt{if} 是条件判断，\texttt{break} 是循环控制语句。

\textbf{知识点：} \texttt{for} 常用于遍历序列或可迭代对象；\texttt{while} 根据条件反复执行。
\end{question}

\begin{question}{6. MPU6050 的特征}
\textbf{点评：} 这题有一个选项它是I2C总线，所以不能与其他传感器合用包错的，但是A选项说MPU6050组成上包括三轴陀螺仪、三轴加速度计，也不知道对不对，少说一个，草台班子。

\textbf{知识点：} MPU6050 是 6 轴运动跟踪设备，包含三轴陀螺仪、三轴加速度计和数字运动处理器，并带有片上温度传感器。它通常通过 I2C 总线与微控制器通信，也可通过辅助 I2C 总线连接其他传感器。
\end{question}

\begin{question}{7. 图像掩码的作用}
\textbf{点评：} 这题 PPT 上标得比较明显，依旧四选

\textbf{知识点：} 掩码可用于提取感兴趣区域、屏蔽无关区域、辅助结构特征提取，以及制作特殊形状图像。
\end{question}

\begin{question}{8. 图像形态学处理}
\textbf{点评：} 依旧全是，四选。

\textbf{知识点：} 开运算通常是先腐蚀后膨胀，适合去除小噪声；闭运算通常是先膨胀后腐蚀，适合填补小孔洞。
\end{question}

\begin{question}{9. GPIO 编码方式}
\textbf{点评：} 终于来了一个 MOOC 原题的内容。

\textbf{知识点：} 常见 GPIO 编码方式包括 WiringPi、BCM 和 BOARD。BOARD 使用物理引脚编号，BCM 使用 Broadcom 芯片 GPIO 编号。
\end{question}

\begin{question}{10. I2C 通信}
\textbf{点评：} 这题是把 MOOC 里关于起始位、终止位、数据线和时钟线的内容杂合在了一起。刷过题以后会轻松很多。

\textbf{知识点：} I2C 主要使用 SDA 数据线和 SCL 时钟线；总线空闲时二者通常为高电平；起始条件为 SCL 高电平时 SDA 下降沿；终止条件为 SCL 高电平时 SDA 上升沿。
\end{question}

\section{三、填空题}

填空题应该是这张卷子的最大痛点

\begin{question}{1. 等待键盘输入函数}
\textbf{点评：} 后面判断题里甚至出现了 \texttt{cv2.waitKey()}，可以直接借来用。

\textbf{知识点：} OpenCV 中常用 \texttt{cv2.waitKey()} 等待键盘输入或控制窗口刷新间隔。
\end{question}

\begin{question}{2. 树莓派开发的 SSH 工具和远程桌面工具}
\textbf{点评：} 当时只记住了远程桌面工具，SSH 工具没有记牢。

\textbf{知识点：} 常见 SSH 工具是 PuTTY，远程桌面工具可使用 VNC Viewer。
\end{question}

\begin{question}{3. 查看 IP 地址的命令}
\textbf{点评：} 这题比较轻松，答案是 \texttt{ipconfig}。

\textbf{知识点：} Windows 下常用 \texttt{ipconfig} 查看网络配置信息；Linux/树莓派环境中也常见 \texttt{ifconfig} 或 \texttt{ip addr}。
\end{question}

\begin{question}{4. 树莓派供电接口与其他接口}
\textbf{点评：} 这题是个草台班子，ppt写的Pype-c，mooc写的Type-c，服了。

\textbf{知识点：} 树莓派 4B 常见供电接口为 USB Type-C；板上还包含 GPIO、HDMI、摄像头接口、显示屏接口、网口、USB 等接口。
\end{question}

\begin{question}{5. GPIO 常态下的高低电平}
\textbf{点评：} 这题不太记得了，好像是写个Pull-up。 如果是上拉输入，常态通常为高电平；如果是下拉输入，常态通常为低电平。

\textbf{知识点：} Pull-up 会把输入端默认拉到高电平，按键或外部电路触发时可能变为低电平；Pull-down 会把输入端默认拉到低电平。
\end{question}

\begin{question}{6. 返回指定 pin 电平状态的函数}
\textbf{点评：}  我猜想考input（），但是没证据。

\textbf{知识点：} \texttt{GPIO.input()} 用于读取某个 GPIO 当前电平状态；\texttt{GPIO.output()} 用于设置某个 GPIO 引脚输出高电平或低电平。
\end{question}

\begin{question}{7. I2C 总线上最多可接几个 PCF8591}
\textbf{点评：} 这题刷过 MOOC 的话会比较熟，答案是 8 个。

\textbf{知识点：} PCF8591 通过地址引脚组合改变地址，常见情况下可在同一 I2C 总线上挂接多个器件。
\end{question}

\begin{question}{8. 有源蜂鸣器与无源蜂鸣器}
\textbf{点评：} ppt原话

\textbf{知识点：} 有源蜂鸣器内置振荡器，通常发出单一提示性报警声；无源蜂鸣器需要频率脉冲驱动，可发出不同频率的声音。
\end{question}

\begin{question}{9. HSV 的概念}
\textbf{点评：} 考前晚上刚好问了一句，第二天就考到，属于意外捡到。

\textbf{知识点：} H 表示色相/色调 Hue；S 表示饱和度 Saturation；V 或 B 表示亮度 Value/Brightness。
\end{question}

\begin{question}{10. Haar 分类器}
\textbf{点评：} 不太记得这个考啥了，反正不会

\textbf{知识点：} 基于 Haar 特征分类器的对象检测技术是一种有效的对象检测方法，通常通过大量正负样本训练分类器，再用于目标检测。
\end{question}

\begin{question}{11. Python 函数定义关键字}
\textbf{点评：} 这题属于基础题，稳住就能拿分。

\textbf{知识点：} Python 使用 \texttt{def} 定义函数。
\end{question}

\begin{question}{12. 轻触开关按下和松开的高低电平}
\textbf{点评：} 这题直接痛失两分。我把高低电平写反了，整张卷子里凡是涉及高低电平的地方好像都错了。

\textbf{知识点：} 在常见上拉输入连接中，按键松开时检测到高电平，按键按下时接地并检测到低电平。
\end{question}

\section{四、大题}

判断题没有单独整理：一是数量太多，二是很多说法只要认真刷 MOOC 就能快速判断,而且答案很绝对，基本上那些只要、仅仅这些话的都是错的。

\subsection{大题一：补全代码}

\begin{question}{OpenCV 目标识别代码补全}
\textbf{问题：} 题面要求根据给出的目标识别流程补全 5 个空，核心围绕图像缩放、色彩空间转换、图像腐蚀、边界查找和质心计算。

\begin{verbatim}
开启摄像头：VideoStream(0).start() 或 cv2.VideoCapture(0)
读取一帧图像：vs.read()
图像尺寸缩放：(1)
色彩空间转换：(2)
指定感兴趣区域 ROI：cv2.inRange(hsv, colorLower, colorUpper)
图像腐蚀：(3)
找到指定区域的边界：(4)
计算轮廓面积并取最大：max(cnts, key=cv2.contourArea)
找到区域质心：(5)
画圆及圆心：cv2.circle()
\end{verbatim}

\textbf{点评：} 前面几个空比较顺，后面两个很折磨了，不知道写哪个，因为绝大部分考前都是背诵，没理解那个函数是干啥的，换个说法就不会了

\textbf{知识点：} 这类流程通常包括 \texttt{VideoStream(0).start()} 或 \texttt{cv2.VideoCapture(0)} 开启摄像头，\texttt{vs.read()} 读取帧，\texttt{cv2.cvtColor()} 转换色彩空间，\texttt{cv2.inRange()} 做颜色阈值化，形态学操作去噪，\texttt{cv2.findContours()} 查找轮廓，\texttt{cv2.moments()} 计算矩并求质心，最后用 \texttt{cv2.circle()} 标记目标。
\end{question}

\subsection{大题二：往年真题}

这一题是 2021-2022 真题的原题，基本靠背诵。

\begin{figure}[H]
    \centering
    \includegraphics[width=0.5\linewidth]{img1.png}
    \label{fig:innovation-old-exam}
\end{figure}

\subsection{大题三：超声波传感器}

\begin{question}{简述超声波传感器的原理以及应用场景}
\textbf{点评：} 无语了，看都没看，只能瞎编。

\textbf{知识点：} 超声波传感器通过发射超声波并接收反射回波，根据传播时间和声速计算距离。常见应用包括避障、测距、液位检测、机器人导航和倒车雷达等。
\end{question}

\subsection{大题四：项目题}

\begin{question}{围绕自己的项目作答}
\textbf{问题：} 题目要求简述项目功能，写出使用的传感器及硬件结构图，并画出流程图。

\textbf{点评：} 这个每年都考，有时候还问你项目遇到过什么bug,怎么解决的，所以一定要记一记，这个还是很ok的。

\end{question}

\subsection{大题五：补全代码}

\begin{question}{双色 LED 与 PWM 代码补全}
\textbf{问题：} 第一问考补全连线图(答题卡给了你图，考你怎么接线)；第二问题面给出双色 LED 的控制程序，要求补全 GPIO 初始化、PWM 对象创建、PWM 启动、颜色循环和资源释放函数中的空。

\textbf{点评：} 双色LED纯背诵题。

\textbf{知识点：} 双色 LED 可通过两个 GPIO 引脚分别控制红色和绿色通道，再通过 PWM 调节亮度。按题面空位整理如下：

\begin{verbatim}
import RPi.GPIO as GPIO
import time

colors = [0xFF00, 0x00FF, 0x0FF0, 0xF00F]
makerobo_pins = (11, 12)  # PIN 物理管脚：11 红、12 绿

# GPIO 基础初始化
______(1)__________      # 使用物理引脚编号
GPIO.setwarnings(False)
______(2)__________      # 设置引脚为输出模式
GPIO.output(makerobo_pins, GPIO.LOW)

# 创建两路 PWM 对象，频率 2KHz
______(3)__________
______(4)__________

# PWM 初始占空比 0（灯熄灭）
______(5)__________
______(6)__________

def makerobo_pwm_map(x, in_min, in_max, out_min, out_max):
    return (x - in_min) * (out_max - out_min) / (in_max - in_min) + out_min

def makerobo_set_Color(col):
    R_val = col >> 8
    G_val = col & 0x00FF
    R_val = makerobo_pwm_map(R_val, 0, 255, 0, 100)
    G_val = makerobo_pwm_map(G_val, 0, 255, 0, 100)
    p_R.ChangeDutyCycle(R_val)
    p_G.ChangeDutyCycle(G_val)

def makerobo_loop():
    while True:
        ______(7)__________:
            ______(8)__________
            ______(9)__________

def makerobo_destroy():
    ______(10)__________
    ______(11)__________
    GPIO.output(makerobo_pins, GPIO.LOW)
    GPIO.cleanup()

if __name__ == "__main__":
    try:
        makerobo_loop()
    except KeyboardInterrupt:
        makerobo_destroy()
\end{verbatim}

如果按常见写法补全，大致是：\texttt{GPIO.setmode(GPIO.BOARD)}，\texttt{GPIO.setup(makerobo\_pins, GPIO.OUT)}，\texttt{p\_R = GPIO.PWM(makerobo\_pins[0], 2000)}，\texttt{p\_G = GPIO.PWM(makerobo\_pins[1], 2000)}，\texttt{p\_R.start(0)}，\texttt{p\_G.start(0)}，\texttt{for col in colors}，\texttt{makerobo\_set\_Color(col)}，\texttt{time.sleep(0.5)}，\texttt{p\_G.stop()}，\texttt{p\_R.stop()}。
\end{question}


总结下来就是，ppt上的细节知识点，尤其是打红的一定要细看，然后大题每年考的都比较固定。